\chapter*{Liste des abréviations}
\addcontentsline{toc}{chapter}{Liste des abréviations}
\begin{longtable}{@{}p{3cm}p{11.5cm}@{}}
\toprule
\textbf{Sigle} & \textbf{Signification}\\
\midrule
\endhead
ANEO & Réseau de capteurs environnementaux de l'écosystème Tukhnanutha\\
API & Interface de programmation (\emph{Application Programming Interface})\\
ATOM Wind & Application web de restitution du vent 3D (agriculture, éolien)\\
DEM & Modèle numérique de terrain (\emph{Digital Elevation Model})\\
DSO & Objectif de performance mesurable (critère d'acceptation du modèle)\\
EMA & Moyenne exponentielle des poids (\emph{Exponential Moving Average})\\
ERA5 & Cinquième réanalyse atmosphérique du CEPMMT (ECMWF)\\
FAO-56 & Méthode Penman-Monteith de la FAO pour l'évapotranspiration\\
FFT & Transformée de Fourier rapide\\
FNO & Opérateur neuronal de Fourier (\emph{Fourier Neural Operator})\\
GEFS & Ensemble de prévision global de la NOAA (31 membres)\\
GFS & Système de prévision global de la NOAA (\emph{Global Forecast System})\\
IDW & Pondération inverse à la distance (\emph{Inverse Distance Weighting})\\
IGRA & Archive mondiale des radiosondages (\emph{Integrated Global Radiosonde Archive})\\
INM & Institut National de la Météorologie (Tunisie)\\
MCP & Protocole de contexte de modèle (\emph{Model Context Protocol})\\
METAR & Message d'observation météorologique d'aérodrome\\
OMM & Organisation Météorologique Mondiale\\
OS & Objectif spécifique du cahier des charges\\
OSSE & Expérience de simulation de systèmes d'observation\\
PC-FNO3D & Opérateur neuronal de Fourier 3D contraint par la physique\\
RMSE & Racine de l'erreur quadratique moyenne\\
SYNOP & Message d'observation synoptique de station au sol\\
\bottomrule
\end{longtable}
